\documentclass[../../main.tex]{subfiles}
\begin{document}

\begin{flushright}
\footnotesize\textit{【自動文字起こし・要確認】}
\end{flushright}

[解] $C: x^2+y^2=1$, $A(1, 0)$ とする。時刻 $t$ での $P, Q, R$ の位置は
\[ P_t(\cos mt, \sin mt), Q_t(\cos t, \sin t), R_t(\cos(-2t), \sin(-2t)) \]
とかける。$\angle PQR = \angle R$ の時、$P_t R_t$ が直径になる。
ゆえ、$n \in \mathbb{Z}$ として
\[ mt = -2t + (2n+1)\pi \implies (m+2)t = (2n+1)\pi \cdots \text{①} \]
又 $A R_t \perp O Q_t$ から、$l \in \mathbb{Z}$ として、
\[ t = -2t + \dfrac{\pi}{2} + l\pi \implies t = \dfrac{\pi}{3}\left(\dfrac{1}{2} + l\right) = \dfrac{\pi}{6}(1+2l) \cdots \text{②} \]
②を①に代入して
\[ \dfrac{\pi}{6}(m+2)(1+2l) = (2n+1)\pi \cdots \text{③} \]
②及び $0 \le t \le 2\pi$ から $l = 0, 1, 2, 3, 4, 5$ である。
\begin{itemize}
\item $1^\circ \ l=0$: ③から $(m+2) = 6(2n+1)$ だから $(m, n) = (4, 0)$
\item $2^\circ \ l=1$: ③から $(m+2) = 2(2n+1)$ だから $(m, n) = (4, 1), (8, 2)$
\item $3^\circ \ l=2$: ③から $5(m+2) = 6(2n+1)$ だから $(m, n) = (4, 2)$
\item $4^\circ \ l=3$: ③から $7(m+2) = 6(2n+1)$ だから $(m, n) = (4, 3)$
\item $5^\circ \ l=4$: ③から $3(m+2) = 2(2n+1)$ だから $(m, n) = (4, 4), (8, 7)$
\item $6^\circ \ l=5$: ③から $11(m+2) = 6(2n+1)$ だから $(m, n) = (4, 5)$
\end{itemize}
以上まとめて、
\[ (t, m) = \left(\dfrac{\pi}{6}, 4\right), \left(\dfrac{1}{2}\pi, 4\right), \left(\dfrac{1}{2}\pi, 8\right), \left(\dfrac{5}{6}\pi, 4\right), \left(\dfrac{7}{6}\pi, 4\right), \left(\dfrac{3}{2}\pi, 4\right), \left(\dfrac{3}{2}\pi, 8\right), \left(\dfrac{11}{6}\pi, 4\right) \]

\end{document}
